% ===== PRÉAMBULE CANONIQUE — à coller VERBATIM en tête de CHAQUE .tex =====
\documentclass[11pt,a4paper]{article}
\usepackage[utf8]{inputenc}\usepackage[T1]{fontenc}\usepackage{lmodern}
\usepackage[french]{babel}
\usepackage{amsmath,amssymb,mathtools}\usepackage{icomma}
\usepackage{siunitx}\sisetup{locale=FR}  % \qty{}{}, \si{}, \ang{}
\usepackage{xcolor}\usepackage{colortbl}
\usepackage{tikz}\usepackage{pgfplots}\pgfplotsset{compat=1.18}
\usepackage[siunitx,europeanresistors]{circuitikz} % circuits (élec.)
\usepackage[most]{tcolorbox}\usepackage{enumitem}\usepackage{array,booktabs}
\usepackage[a4paper,margin=2.2cm]{geometry}
\usetikzlibrary{arrows.meta,calc,angles,quotes,positioning,patterns,%
  backgrounds,shapes.geometric,decorations.pathmorphing,decorations.markings,intersections}
\definecolor{primary}{RGB}{222,49,99}\definecolor{primarydark}{RGB}{150,22,55}
\definecolor{defcol}{RGB}{0,114,178}\definecolor{methcol}{RGB}{52,92,148}
\definecolor{excol}{RGB}{0,135,90}\definecolor{attncol}{RGB}{200,40,55}
\tcbset{boxbase/.style={enhanced,breakable,boxrule=0.9pt,arc=2.5pt,
  left=3mm,right=3mm,top=2.3mm,bottom=2.3mm,fonttitle=\bfseries,coltitle=white}}
% --- boîtes TUTORIEL (noms EXACTS, ne pas renommer) ---
\newtcolorbox{cle}[1][]{boxbase,colback=defcol!6,colframe=defcol,
  title={\textsc{À retenir}\ifx\relax#1\relax\relax\else\textmd{ : \itshape #1}\fi}}
\newtcolorbox{methode}[1][]{boxbase,colback=methcol!7,colframe=methcol,sharp corners,
  title={\textsc{Méthode}\ifx\relax#1\relax\relax\else\textmd{ --- \itshape #1}\fi}}
\newtcolorbox{exemplebox}[1][]{boxbase,colback=excol!5,colframe=excol,
  title={\textsc{Exemple}\ifx\relax#1\relax\relax\else\textmd{ : \itshape #1}\fi}}
\newtcolorbox{piege}[1][]{boxbase,colback=attncol!6,colframe=attncol,
  title={\textsc{Piège}\ifx\relax#1\relax\relax\else\textmd{ : \itshape #1}\fi}}
\newtcolorbox{remarque}[1][]{boxbase,colback=black!4,colframe=black!45,
  title={\textsc{Remarque}\ifx\relax#1\relax\relax\else\textmd{ : \itshape #1}\fi}}
% --- boîtes EXERCICE / CORRIGÉ (noms EXACTS, auto-comptées) ---
\newtcolorbox[auto counter]{exo}[1][]{boxbase,colback=primary!4,colframe=primary,
  title={\textsc{Exercice}~\thetcbcounter\ifx\relax#1\relax\relax\else\textmd{ --- \itshape #1}\fi}}
\newtcolorbox[auto counter]{corrige}[1][]{boxbase,colback=excol!4,colframe=excol,
  title={\textsc{Corrigé}~\thetcbcounter\ifx\relax#1\relax\relax\else\textmd{ --- \itshape #1}\fi}}
\newcommand{\vv}[1]{\vec{#1}}\newcommand{\ds}{\displaystyle}
\newcommand{\R}{\mathbb{R}}\newcommand{\N}{\mathbb{N}}\newcommand{\Z}{\mathbb{Z}}
\begin{document}
\sisetup{per-mode=symbol}

\begin{exo}[objet placé en $2F$]
Un objet réel est placé exactement en $2F$ (à la distance $2f$) devant une lentille convergente.
\begin{enumerate}[label=\alph*)]
  \item En traçant les rayons particuliers, en quel point de l'axe se forme l'image ?
  \item Précise sa nature (réelle/virtuelle), son sens et sa taille par rapport à l'objet.
\end{enumerate}
\end{exo}

\begin{exo}[quel rayon émergent est correct ?]
Le rayon incident représenté passe par le \textbf{foyer objet $F$} avant d'atteindre la lentille
convergente. Lequel des quatre tracés proposés après la lentille est correct ?
\begin{center}
\begin{tikzpicture}[scale=0.85,>=Stealth,font=\footnotesize]
  \draw[black!35,dash pattern=on 3pt off 2pt] (-4.6,0)--(4.2,0);
  \draw[line width=1.5pt,<->,black!65] (0,-2)--(0,2);
  \fill[black] (0,0) circle (1.2pt); \node[below right] at (0,0) {$O$};
  \fill[primary] (2,0) circle (1.5pt); \node[below] at (2,0) {$F'$};
  \fill[primary] (-2,0) circle (1.5pt); \node[below] at (-2,0) {$F$};
  % incident passant par F, atteignant la lentille en (0,1)
  \draw[black!80,line width=1.0pt,->] (-4.2,-1.1)--(0,1);
  \node[black!80,left] at (-4.2,-1.1) {incident};
  % candidates depuis (0,1)
  \draw[defcol,line width=0.8pt,->] (0,1)--(3.6,1) node[right]{1};
  \draw[defcol,line width=0.8pt,->] (0,1)--(3.4,-0.7) node[right]{2};
  \draw[defcol,line width=0.8pt,->] (0,1)--(2.4,2.26) node[right]{3};
  \draw[defcol,line width=0.8pt,->] (0,1)--(2.2,-1.6) node[right]{4};
\end{tikzpicture}
\end{center}
\end{exo}

\begin{exo}[compléter le tableau des natures]
Pour une lentille \textbf{convergente}, indique la nature de l'image (réelle/virtuelle, droite/renversée,
agrandie/réduite/même taille) selon la position de l'objet réel :
\begin{enumerate}[label=\alph*)]
  \item objet à l'infini ;
  \item objet au-delà de $2F$ ;
  \item objet entre $F$ et $2F$ ;
  \item objet entre $O$ et $F$.
\end{enumerate}
\end{exo}

\begin{exo}[construction avec une divergente]
Un objet réel $AB$ est placé devant une lentille \textbf{divergente}, à une distance supérieure à
$|f|$.
\begin{enumerate}[label=\alph*)]
  \item À l'aide du rayon par $O$ et du rayon parallèle (qui semble venir de $F'$), indique de quel
        côté de la lentille se forme l'image.
  \item Donne sa nature complète (réelle/virtuelle, sens, taille).
\end{enumerate}
\end{exo}

\begin{exo}[l'objet s'approche de la lentille]
On déplace un objet réel depuis l'infini jusqu'au centre optique d'une lentille convergente. Décris
l'évolution de l'image :
\begin{enumerate}[label=\alph*)]
  \item tant que l'objet reste au-delà du foyer $F$ (de l'infini jusqu'à $F$) ;
  \item une fois l'objet passé entre $F$ et $O$.
\end{enumerate}
\end{exo}

\end{document}
